\documentclass[10pt,a4paper]{article}
\usepackage[top=1.8cm,bottom=1.8cm,left=1.5cm,right=1.5cm]{geometry}
\usepackage[dvipsnames,svgnames,x11names]{xcolor}
\usepackage{amsmath,amssymb,mathtools}
\usepackage{multicol}
\usepackage{pifont}
\usepackage{hyperref}

\setlength{\parindent}{0pt}
\setlength{\parskip}{3pt plus 1pt minus 1pt}
\setlength{\columnsep}{1.2cm}

\definecolor{mainbrown}{RGB}{128, 40, 20}
\definecolor{maingreen}{RGB}{85, 107, 47}

\hypersetup{
    pdfauthor={CPES 2},
    pdftitle={FICHE DE COURS N°1 - LES SERIES NUMERIQUES},
    hidelinks
}

\begin{document}

% ==================== PAGE 1 & 2 ====================
\begin{multicols}{2}

\begin{flushright}
\footnotesize CPES 2
\end{flushright}

\vspace{-0.2cm}
\begin{center}
\fbox{\parbox{0.95\linewidth}{\centering \bfseries\large \textcolor{mainbrown}{FICHE DE COURS N\textdegree 1 -- LES SERIES NUMERIQUES}}}
\end{center}
\vspace{0.1cm}

Dans ce chapitre, nous considérons des suites ($u_n$) à valeurs dans $\mathbb{R}$.

\vspace{0.2cm}
\textbf{\textcolor{mainbrown}{I.\quad SÉRIES NUMÉRIQUES}}

\vspace{0.1cm}
\textbf{\textcolor{maingreen}{A. DÉFINITIONS}}

\textbf{\textit{Définition 1 :}} Étant donnée une suite $(u_n)_{n \ge n_0}$, on appelle série de terme général $u_n$, notée $\sum_{n \ge n_0} u_n$ ou $\sum u_n$ lorsque $n_0 = 0$, la suite des sommes partielles $(S_n)_{n \ge n_0}$ où pour tout entier naturel $n \ge n_0$, on a : $S_n = \sum_{k=n_0}^n u_k$.

\textbf{\textit{Exemple 1 :}} Soit ($u_n$) la suite définie pour tout $n \ge 1$ par $u_n = \frac{1}{n}$.

La série de terme général $u_n$ est notée $\sum_{n \ge 1} \frac{1}{n}$, et est appelée la série harmonique.

Les premières sommes partielles sont :
\begin{equation*}
S_1 = 1,\; S_2 = 1 + \frac{1}{2} = \frac{3}{2},\; S_3 = 1 + \frac{1}{2} + \frac{1}{3} = \frac{11}{6},\; \dots
\end{equation*}

\textbf{\textit{Définition 2 :}} On dit que $\sum u_n$ converge si la suite ($S_n$) admet une limite finie dans $\mathbb{R}$.

Dans ce cas, la limite de la suite ($S_n$) est appelée somme de la série, et est notée $\sum_{k=0}^{+\infty} u_k$.

On a ainsi $\sum_{k=0}^{+\infty} u_k = \lim_{n \to +\infty} \sum_{k=0}^n u_k$.

Dans le cas contraire, on dit que la série diverge.

\textbf{\textit{Exemple 2 :}} Soit ($u_n$) la suite définie pour tout entier naturel $n$ par $u_n = 0,5^n$.

Alors $\forall n \in \mathbb{N},\; S_n = \sum_{k=0}^n 0,5^k = \frac{1-0,5^{n+1}}{1-0,5} = 2(1 - 0,5^{n+1})$.

Puisque $-1 < 0,5 < 1$, on a $\lim_{n \to +\infty} 0,5^{n+1} = 0$, donc la série $\sum_{n \ge 0} 0,5^n$ converge vers 2 car $\lim_{n \to +\infty} S_n = \sum_{k=0}^{+\infty} 0,5^k = 2$.

\textbf{\textit{Remarque 1 :}} Étudier la nature d'une série, c'est déterminer si la série converge ou diverge.

\textbf{\textit{Attention :}} Les sommes infinies ne se manipulent pas comme les sommes finies (puisqu'en réalité, ce sont des limites, et il faut donc toujours s'assurer de la convergence). En effet, l'écriture $\sum_{k=0}^{+\infty} u_k$ n'a de sens que si la série converge, alors que l'écriture $\sum u_n$ a bien un sens, puisqu'elle désigne une suite.

C'est pourquoi on calculera (presque) toujours les sommes partielles, qui sont des sommes finies, avant de passer à la limite.

\textbf{\textit{Exemple 3 :}} Montrer que la série harmonique, de terme général $\frac{1}{n}$ est divergente.

Pour tout $n \ge 1$, notons $H_n = \sum_{k=1}^n \frac{1}{k}$, et soit $k \in \mathbb{N}^*$ :

\vfill
\centerline{\footnotesize 1}

\newpage

\begin{flushright}
\footnotesize CPES 2
\end{flushright}
\vspace{-0.2cm}

\begin{equation*}
\ln(k+1) - \ln(k) = \ln\left(\frac{k+1}{k}\right) = \ln\left(1 + \frac{1}{k}\right) \le \frac{1}{k}
\end{equation*}
Ainsi, pour tout $k \ge 1$, $\ln(k+1) - \ln(k) \le \frac{1}{k}$

En additionnant ces inégalités membre à membre, on obtient :
\begin{equation*}
\sum_{k=1}^n \ln(k+1) - \ln(k) \le \sum_{k=1}^n \frac{1}{k}
\end{equation*}
La somme de gauche étant une somme télescopique, on a :
\begin{equation*}
\sum_{k=1}^n \ln(k+1) - \ln(k) = \ln(n+1) - \ln 1 = \ln(n+1)
\end{equation*}
Ainsi, pour tout $n \ge 1$, $H_n \ge \ln(n+1)$.

Puisque $\lim_{n \to +\infty} \ln(n+1) = +\infty$, par comparaison on a $\lim_{n \to +\infty} \sum_{k=1}^n \frac{1}{k} = +\infty$.

\vspace{0.1cm}
\textbf{\textcolor{maingreen}{B. PROPRIÉTÉS}}

\textbf{\textit{Proposition 1 :}} 1) On ne change pas la nature d'une série (convergente ou divergente) en supprimant un nombre fini de termes.

2) Si la série $\sum_{n \ge n_0} u_n$ converge alors $\lim_{n \to +\infty} u_n = 0$.

\textbf{\textit{Démonstration partielle :}} 2) Pour tout $n$, notons $S_n = \sum_{k=0}^n u_k$.

Alors, pour tout $n \ge 1$, on a $u_n = S_n - S_{n-1}$.

Si la série $\sum u_n$ converge, alors la suite ($S_n$) admet, par définition, une limite que l'on note $S$. Mais alors $\lim_{n \to +\infty} S_n = \lim_{n \to +\infty} S_{n-1} = S$, et donc $\lim_{n \to +\infty} u_n = S - S = 0$.

\textbf{\textit{Attention :}} La réciproque de cette proposition est fausse.

En effet, on a vu que la série harmonique diverge, alors que $\lim_{n \to +\infty} \frac{1}{n} = 0$.

\textbf{\textit{Remarque 2 :}} La contraposée de cette proposition dit que si $\lim_{n \to +\infty} u_n \ne 0$, alors la série $\sum_{n \ge n_0} u_n$ diverge.

\textbf{\textit{Exemple 4 :}} La série $\sum \frac{2n}{n+1}$ diverge car $\lim_{n \to +\infty} \frac{2n}{n+1} = 2 \ne 0$.

\textbf{\textit{Proposition 2 (admise) :}} 1) Si les séries $\sum_{n \ge n_0} u_n$ et $\sum_{n \ge n_0} v_n$ convergent alors $\sum_{n \ge n_0} (u_n + v_n)$ converge et dans ce cas : $\sum_{n \ge n_0} (u_n + v_n) = \sum_{n \ge n_0} u_n + \sum_{n \ge n_0} v_n$.

2) Soit $\lambda \in \mathbb{R}^*$ les séries $\sum_{n \ge n_0} u_n$ et $\sum_{n \ge n_0} \lambda u_n$ sont de même nature (convergente ou divergente), et si la série $\sum_{n \ge n_0} u_n$ converge alors $\sum_{n=n_0}^{+\infty} \lambda u_n = \lambda \sum_{n=n_0}^{+\infty} u_n$.

\textbf{\textit{Attention :}} La réciproque du $1^{\text{er}}$ point est fausse. En effet, les séries $\sum_{n \ge 1} \frac{1}{n}$ et $\sum_{n \ge 1} -\frac{1}{n}$ sont toutes deux divergentes alors que la série $\sum_{n \ge 1} \frac{1}{n} + \left(-\frac{1}{n}\right)$ converge vers 0.

\vspace{0.1cm}
\textbf{\textcolor{maingreen}{C. SÉRIES À TERMES POSITIFS}}

Dans toute cette partie ($u_n$) et ($v_n$) désignent deux suites à termes positifs

\vfill
\centerline{\footnotesize 2}

\end{multicols}

\newpage

% ==================== PAGE 3 & 4 ====================
\begin{multicols}{2}

\begin{flushright}
\footnotesize CPES 2
\end{flushright}
\vspace{-0.2cm}

\textbf{\textit{Proposition 3 :}} La série $\sum_{n \ge n_0} u_n$ est convergente, si et seulement si, la suite des sommes partielles $(S_n)$ est majorée.

\textbf{\textit{Démonstration :}} Notons $S_n = \sum_{k=0}^n u_k$.

On a alors $S_{n+1} - S_n = u_{n+1}$ et $u_{n+1} > 0$. La suite $(S_n)$ est donc croissante.

D'après les théorèmes sur les suites monotones, $(S_n)$ converge si et seulement si la suite $(S_n)$ est majorée.

\textbf{\textit{Remarque 3 :}} Pour majorer $(S_n)$, on commence en général par majorer $(u_n)$. On somme alors ces majorants pour en déduire un majorant de $(S_n)$. On dispose ainsi du théorème suivant :

\textbf{\textit{Proposition 4 (théorèmes de comparaison) :}} Si pour tout $n \ge n_0$, on a $0 \le u_n \le v_n$, alors,
\begin{itemize}
\item[\checkmark] si la série $\sum v_n$ converge, la série $\sum u_n$ est également convergente.

Dans ce cas, $\sum_{k=n_0}^{+\infty} u_k \le \sum_{k=n_0}^{+\infty} v_k$.
\item[\checkmark] si la série $\sum u_n$ diverge, alors la série $\sum v_n$ est également divergente.
\end{itemize}

\textbf{\textit{Démonstration :}} Si on note $S_n = \sum_{k=n_0}^n u_k$ et $T_n = \sum_{k=n_0}^n v_k$, on a, pour tout $n \ge n_0$ :

$S_n \le T_n$ (addition des inégalités). De plus, la suite $(T_n)$ est également croissante, de limite T. Donc pour tout $n \ge n_0$, $T_n \le T$.

Donc $\forall n \ge n_0$, $S_n \le T_n \le T$. La suite $(S_n)$ est donc majorée, et d'après le théorème précédent, la série $\sum u_n$ converge. L'inégalité précédente donne alors $\sum_{n=n_0}^{+\infty} u_n = \lim_{n \to +\infty} S_n \le T = \lim_{n \to +\infty} T_n = \sum_{n=n_0}^{+\infty} v_n$.

\textbf{\textit{Proposition 5 (admise) :}} Si $u_n \underset{+\infty}{\sim} v_n$ alors $\sum_{n \ge n_0} u_n$ et $\sum_{n \ge n_0} v_n$ sont de même nature.

\textbf{\textit{Exemple 5 :}} La série $\sum_{n \ge 1} \frac{n+1}{3n^2+n+1}$ est à termes positifs. De plus $\frac{n+1}{3n^2+n+1} \underset{+\infty}{\sim} \frac{1}{3n}$ et $\sum_{n \ge 1} \frac{1}{3n}$ est une série divergente donc la série $\sum_{n \ge 1} \frac{n+1}{3n^2+n+1}$ diverge.

\textbf{\textit{Proposition 6 (Règle de d'Alembert) :}} Soit $(u_n)$ une suite réelle tel que $u_n > 0$ pour tout entier naturel $n$. On suppose $\lim_{n \to +\infty} \frac{u_{n+1}}{u_n} = \ell$. On a alors les assertions suivantes :
\begin{itemize}
\item[\tiny$\blacksquare$] Si $\ell < 1$, la série de terme général $u_n$ est convergente;
\item[\tiny$\blacksquare$] Si $\ell > 1$, la série de terme général $u_n$ est divergente;
\item[\tiny$\blacksquare$] si $\ell = 1$, nous ne pouvons pas conclure.
\end{itemize}

\textbf{\textit{Démonstration :}} $\forall \varepsilon > 0, \exists N \in \mathbb{N}$ tel que $\forall n \ge N$, $\left|\frac{u_{n+1}}{u_n} - \ell\right| < \varepsilon$.

On en déduit que $\forall \varepsilon > 0, \exists N \in \mathbb{N}$ tel que $\forall n \ge N$, $(\ell - \varepsilon)u_n < u_{n+1} < (\ell + \varepsilon)u_n$.

On en déduit encore que :

$\forall \varepsilon > 0, \exists N \in \mathbb{N}$ tel que $\forall n \ge N$, $(\ell - \varepsilon)^{n-N}u_N < u_n < (\ell + \varepsilon)^{n-N}u_N$.

\vfill
\centerline{\footnotesize 3}

\newpage

\begin{flushright}
\footnotesize CPES 2
\end{flushright}
\vspace{-0.2cm}

Supposons maintenant que $\ell < 1$, et soit $\varepsilon = \frac{1-\ell}{2}$. Alors $\varepsilon > 0$, il existe donc $N \in \mathbb{N}$ tel que $\forall n \ge N$, $u_n < (\ell + \varepsilon)^{n-N}u_N$. Or la série de terme général $(\ell + \varepsilon)^{n-N}u_N = \left(\frac{\ell+1}{2}\right)^{n-N}u_N$ est convergente (on reconnaît le terme général d'une série géométrique convergente), donc la série de terme général $u_n$ est convergente.

Si $\ell > 1$, on raisonne de même avec $\varepsilon = \frac{\ell-1}{2} > 0$, et on a $\exists N \in \mathbb{N}$ tel que $\forall n \ge N$,

$\left(\frac{\ell+1}{2}\right)^{n-N}u_N < u_n$. Or la série de terme général $\left(\frac{\ell+1}{2}\right)^{n-N}u_N$ est divergente (on reconnaît le terme général d'une série géométrique divergente), donc la série de terme général $u_n$ est divergente.

\textbf{\textit{Exemple 6 :}} $\sum \frac{2^n n!}{n^n}$ est une série à termes positifs, et $\frac{u_n}{u_{n+1}} = \frac{2^n n!}{n^n} \times \frac{(n+1)^{n+1}}{2^{n+1}(n+1)!} = \frac{1}{2}\left(\frac{n+1}{n}\right)^n$.

Or $\lim_{n \to +\infty} n \ln\left(1 + \frac{1}{n}\right) = 1$, donc $\lim_{n \to +\infty}\left(\frac{u_n}{u_{n+1}}\right) = \frac{e}{2}$ et donc $\lim_{n \to +\infty}\left(\frac{u_{n+1}}{u_n}\right) = \frac{2}{e} < 1$.

Ce qui prouve que la série converge et donc $\lim_{n \to +\infty}\left(\frac{2^n n!}{n^n}\right) = 0$

\textbf{\textit{Proposition 7 (Théorème de Fubini) :}} Soit $(u_{ij})_{(i,j) \in \mathbb{N}^2}$ une suite double positive.

On suppose que $\forall i \in \mathbb{N}, \sum_{j \in \mathbb{N}} u_{ij}$ converge et $\sum_{i \in \mathbb{N}} \left(\sum_{j=0}^{+\infty} u_{ij}\right)$ converge. On a alors :

$\forall j \in \mathbb{N}, \sum_{i \in \mathbb{N}} u_{ij}$ converge, $\sum_{j \in \mathbb{N}} \left(\sum_{i=0}^{+\infty} u_{ij}\right)$ converge et $\sum_{i=0}^{+\infty} \sum_{j=0}^{+\infty} u_{ij} = \sum_{j=0}^{+\infty} \sum_{i=0}^{+\infty} u_{ij}$.

\textbf{\textit{Remarque 4 :}} Ce théorème signifie que l'on peut échanger deux limites, ce qui n'est pas possible en règle générale.

\textbf{\textit{Exemple 7 :}} On pose pour tout $(i,j) \in \mathbb{N}^2$, $a_{ij} = \begin{cases} 1 & \text{si } i = j \\ -1 & \text{si } j = i + 1 \\ 0 & \text{sinon} \end{cases}$

$\forall i \in \mathbb{N}, \sum_{j \in \mathbb{N}} a_{ij}$ converge et vaut 0, donc $\sum_{i \in \mathbb{N}} \left(\sum_{j=0}^{+\infty} a_{ij}\right)$ converge et vaut 0.

\vspace{0.1cm}
\textbf{\textcolor{maingreen}{D. SÉRIES ALTERNÉES}}

\textbf{\textit{Proposition 8 :}} Soit $(u_n)$ une suite réelle décroissante convergeant vers 0. Alors la série $\sum (-1)^n u_n$ est convergente.

\textbf{\textit{Démonstration :}} Soit $n \in \mathbb{N}$.

Les relations $S_{2n+2} = S_{2n} - u_{2n+1} + u_{2n+2}$ et $S_{2n+3} = S_{2n+1} + u_{2n+2} - u_{2n+3}$, ainsi que la décroissance de la suite $(u_n)$ entraînent : $S_{2n+2} \le S_{2n}$ et $S_{2n+1} \le S_{2n+3}$.

La suite $(S_{2n+1})$ est donc croissante et $(S_{2n})$ est décroissante. Par ailleurs, on a :

$\forall n \in \mathbb{N}, S_{2n} - S_{2n+1} = u_{2n+1}$, donc $\lim_{n \to +\infty} (S_{2n} - S_{2n+1}) = 0$.

Les suites $(S_{2n})$ et $(S_{2n+1})$ sont ainsi adjacentes, donc convergent vers un même réel S. Cela prouve la convergence de la suite $(S_{2n})$ et donc de la série $\sum (-1)^n u_n$.

\textbf{\textit{Exemple 8 :}} Pour tout $\alpha > 0$, la série $\sum_{n \ge 1} \frac{(-1)^{n-1}}{n^\alpha}$ converge.

\vfill
\centerline{\footnotesize 4}

\end{multicols}

\newpage

% ==================== PAGE 5 & 6 ====================
\begin{multicols}{2}

\begin{flushright}
\footnotesize CPES 2
\end{flushright}
\vspace{-0.2cm}

\textbf{\textcolor{mainbrown}{II.\quad SERIES DE REFERENCE}}

\vspace{0.1cm}
\textbf{\textcolor{maingreen}{A. SÉRIES GÉOMÉTRIQUES ET DÉRIVÉES}}

\textbf{\textit{Définition 3 :}} Pour tout entier $q$, la série $\sum q^n$ s'appelle série géométrique de raison $q$.

\textbf{\textit{Proposition 9 :}} La série $\sum q^n$ est convergente si et seulement si $|q| < 1$.

Dans ce cas, $\sum_{n=0}^{+\infty} q^n = \frac{1}{1-q}$.

Plus généralement, la série $\sum_{n \ge p} q^n$ est convergente si et seulement si $|q| < 1$, et dans ce cas, $\sum_{n=p}^{+\infty} q^n = \frac{q^p}{1-q}$.

\textbf{\textit{Démonstration :}} La suite $(q^n)$ converge vers 0 si et seulement si $|q| < 1$.

Par condition nécessaire de convergence, la série $\sum q^n$ ne peut pas converger si $|q| > 1$.

Supposons alors que $|q| < 1$. Notons $S_n = \sum_{k=0}^n q^k$. On a $S_n = \frac{1-q^{n+1}}{1-q}$. Or, $\lim_{n \to +\infty} q^{n+1} = 0$ car $|q| < 1$, donc la suite $(S_n)$ converge vers $\frac{1}{1-q}$ : la série converge, et sa somme vaut $\frac{1}{1-q}$.

\textbf{\textit{Proposition 10 :}} 1) La série $\sum_{n \ge 1} n q^{n-1}$ converge si et seulement si $|q| < 1$. Dans ce cas, $\sum_{n=1}^{+\infty} n q^{n-1} = \frac{1}{(1-q)^2}$ (série géométrique dérivée première).

2) La série $\sum_{n \ge 2} n(n-1) q^{n-2}$ converge si et seulement si $|q| < 1$. Dans ce cas, $\sum_{n=2}^{+\infty} n(n-1) q^{n-2} = \frac{2}{(1-q)^3}$ (série géométrique dérivée deuxième).

\textbf{\textit{Démonstration :}} 1) Si $|q| > 1$, la suite $(n q^{n-1})$ ne converge pas vers 0, donc la série $\sum_{n \ge 1} n q^{n-1}$ ne peut converger.

Pour tout $x \in ]-1; 1[$, notons $T_n(x) = \sum_{k=0}^n x^k$.

$T_n$ est une fonction dérivable sur $]-1; 1[$, et on a $T_n'(x) = \sum_{k=1}^n k x^{k-1}$.

D'autre part, $T_n(x) = \frac{1-x^{n+1}}{1-x}$.

On a donc également $T_n'(x) = \frac{-(n+1)x^n(1-x)+(1-x^{n+1})}{(1-x)^2} = \frac{1-(n+1)x^n+n x^{n+1}}{(1-x)^2}$.

$\lim_{n \to +\infty} |n x^n| = \lim_{n \to +\infty} n |x|^n$
\begin{align*}
&= \lim_{n \to +\infty} n e^{n \ln |x|} \\
&= \frac{1}{\ln |x|} \lim_{X \to -\infty} X e^X
\end{align*}
En posant $X = n \ln |x|$ ($\lim_{n \to +\infty} X = -\infty$ car $\ln |x| < 0$)

Or $\lim_{X \to -\infty} X e^X = 0$, donc $\lim_{n \to +\infty} |n x^n| = 0$.

On en déduit que $\lim_{n \to +\infty} n x^n = 0$.

Ce qui prouve que $\lim_{n \to +\infty} (n+1) x^n = \lim_{n \to +\infty} n x^n + \lim_{n \to +\infty} x^n = 0$ et que :

$\lim_{n \to +\infty} n x^{n+1} = x \lim_{n \to +\infty} n x^n = 0$.

On en déduit que $\lim_{n \to +\infty} T_n'(x) = \frac{1}{(1-x)^2}$.

Ainsi, la série $\sum_{n \ge 1} n x^{n-1}$ converge, et on a bien $\sum_{n \ge 1} n x^{n-1} = \frac{1}{(1-x)^2}$.

\vfill
\centerline{\footnotesize 5}

\newpage

\begin{flushright}
\footnotesize CPES 2
\end{flushright}
\vspace{-0.2cm}

\textbf{\textit{Remarque 5 :}} On remarque que $n q^n = q n q^{n-1}$ et que $n(n-1) q^n = q^2 n(n-1) q^{n-2}$.

Ainsi, si $|q| < 1$, les séries $\sum n q^n$ et $\sum n(n-1) q^n$ convergent également, et :

$\sum_{n=1}^{+\infty} n q^n = \frac{q}{(1-q)^2}$ et $\sum_{n=1}^{+\infty} n(n-1) q^n = \frac{2 q^2}{(1-q)^3}$.

\vspace{0.1cm}
\textbf{\textcolor{maingreen}{B. SÉRIES DE RIEMANN}}

\textbf{\textit{Définition 4 :}} La série de terme général $\frac{1}{n^\alpha}$ ($\alpha \in \mathbb{R}$, $n \in \mathbb{N}^*$) est appelée série de Riemann.

\textbf{\textit{Proposition 11 (admise) :}} La série de Riemann $\sum_{n \ge 1} \frac{1}{n^\alpha}$ converge si et seulement si $\alpha > 1$.

\textbf{\textit{Exemple 9 :}} La série $\sum_{n \ge 1} \frac{1}{n^2}$ est une série de Riemann convergente.

\vspace{0.1cm}
\textbf{\textcolor{maingreen}{C. SÉRIE EXPONENTIELLE}}

\textbf{\textit{Définition 5 :}} La série de terme général $\frac{x^n}{n!}$ ($x \in \mathbb{R}$, $n \in \mathbb{N}$) est appelée série exponentielle.

\textbf{\textit{Proposition 12 (admise) :}} Pour tout réel $x$, la série exponentielle $\sum \frac{x^n}{n!}$ converge, et on a : $\sum_{n=0}^{+\infty} \frac{x^n}{n!} = e^x$.

\textbf{\textit{Méthode :}} Pour déterminer si une série converge ou non, et éventuellement calculer sa limite, on essaiera de la comparer à une des séries usuelles (géométriques, Riemann ou exponentielle).

\textbf{\textit{Exemple 10 :}} Déterminer la nature de la série de terme générale $u_n = \frac{(-3)^{n+1}}{n!}$.

Remarquons tout d'abord que $u_n = -3 \frac{(-3)^n}{n!}$.

La série $\sum_{n \ge 0} \frac{(-3)^n}{n!}$ converge puisqu'il s'agit de la série exponentielle, et sa somme vaut $e^{-3}$.

La série $\sum_{n \ge 0} u_n$ converge donc, et $\sum_{n=0}^{+\infty} \frac{(-3)^{n+1}}{n!} = -3 e^{-3}$.

\textbf{\textit{Remarque 6 :}} Par décalage d'indice, on a également, pour tout réel $x$, $\sum_{n \ge 1} \frac{x^{n-1}}{(n-1)!}$ converge et,

$\sum_{n=1}^{+\infty} \frac{x^{n-1}}{(n-1)!} = e^x$.

\vspace{0.2cm}
\textbf{\textcolor{mainbrown}{III.\quad SERIES ABSOLUMENT CONVERGENTES}}

\vspace{0.1cm}
\textbf{\textit{Définition 6 :}} On dit que la série $\sum u_n$ est absolument convergente, ou converge absolument, si la série à terme reels positifs $\sum |u_n|$ est convergente.

\textbf{\textit{Proposition 13 :}} Toute série absolument convergente est convergente.

\vfill
\centerline{\footnotesize 6}

\end{multicols}

% Page 4 continuation
\hfill{\footnotesize CPES 2}
\vspace{0.2cm}

\noindent
\begin{minipage}[t]{0.48\textwidth}
\textit{\textbf{Démonstration :}} Soit $(u_n)$ telle que la série $\sum |u_n|$ converge.\\[0.15cm]
Pour tout $n \in \mathbb{N}$, les réels $u_n^+ = \frac{|u_n| + u_n}{2}$ et $u_n^- = \frac{|u_n| - u_n}{2}$ sont positifs et vérifient $u_n = u_n^+ - u_n^-$ et $|u_n| = u_n^+ + u_n^-$.\\[0.15cm]
En particulier, $0 \le u_n^+ \le |u_n|$ et $0 \le u_n^- \le |u_n|$.\\[0.15cm]
La convergence de la série de terme général $|u_n|$ entraîne donc la convergence des séries à termes positifs $\sum u_n^+$ et $\sum u_n^-$.\\[0.15cm]
Par suite, $\sum u_n^+ - u_n^-$ est une série convergente, en tant que différence de séries convergentes, donc $\sum u_n$ converge.\\[0.3cm]
\textit{\textbf{Attention :}} La réciproque est fausse. La série $\sum_{n \ge 1} \frac{(-1)^n}{n}$ converge alors qu'elle ne converge pas absolument.\\[0.3cm]
\textit{\textbf{Proposition 14 (admise) :}} Si $u_n = o(v_n)$, alors la convergence absolue de $\sum v_n$ entraîne celle de $\sum u_n$.\\[0.3cm]
\textit{\textbf{Exemple 11 :}} La série $\sum_{n \ge 1} \frac{\ln n}{n 2^n}$ est convergente.\\[0.15cm]
En effet, comme par croissance comparée, $\lim_{n \to +\infty} \frac{\ln n}{n} = 0$, donc $\frac{\ln n}{n 2^n} = o\left(\frac{1}{2^n}\right)$.\\[0.15cm]
Comme la série géométrique, à termes positifs, est convergente, la série $\sum_{n \ge 1} \frac{\ln n}{n 2^n}$ converge (absolument).
\end{minipage}%
\hfill
\begin{minipage}[t]{0.48\textwidth}
\vspace{0pt}
\framebox[\textwidth][l]{%
\begin{minipage}{0.96\textwidth}
\vspace{0.15cm}
If $f$ is \textcolor{CadetBlue}{\underline{asymptotically equivalent}} to \textbf{$g$ for $x$ tending to $c$}, we use the notation $f \sim g$, $x \to c$.\\[0.3cm]
If $f$ is \textbf{negligible with respect to $g$ when $x$ goes to $c$}, the symbol $f = o(g)$, $x \to c$, is used, spoken « \textbf{$f$ is} \textcolor{CadetBlue}{\underline{little o}} \textbf{of $g$ for $x$ tending to $c$} »\\[0.3cm]
$\{u_n\}$ ($u_n$ is in braces) represents the \textcolor{CadetBlue}{\underline{sequence}} that can be generated by using $u_n$ as the $n$th term.\\[0.3cm]
The sum $S = \sum_{i=1}^{+\infty} a_i$ is an \textcolor{CadetBlue}{\underline{infinite series}}. Its value, if one exists, is the limit of the sequence of \textbf{partial sums} $\{s_n\}$ with $s_n = \sum_{i=1}^n a_i$. In that case, the series is said to be convergent.\\
If not, the series is divergent.\\[0.3cm]
For example, the series $\sum_{n \ge 1} a r^{n-1}$, where $a$ and $r$ are constant, is a \textbf{geometric series}.\\[0.15cm]
The series $\sum_{n \ge 1} \frac{1}{n^p}$, where $p$ is a constant, are called \textbf{$p$ series}.\\[0.15cm]
The series whith $p = 1$ is called \textbf{harmonic series}.
\vspace{0.15cm}
\end{minipage}%
}
\end{minipage}

\vspace{0.8cm}

\noindent
\begin{minipage}[t]{0.48\textwidth}
\framebox[\textwidth][l]{%
\begin{minipage}{0.96\textwidth}
\vspace{0.15cm}
The sum $a_1 + a_2 + \dots + a_n$ can be written more compactly using \textcolor{CadetBlue}{\underline{sigma notation}}.\\
We write it as $\sum_{r=1}^n a_r$ (read ``the sum for $r$ equals $1$ up to $n$ of $a_r$'' or ``the sum of all the terms $a_r$ where $r$ takes the values from $1$ to $n$'' or ``the sum of $a_r$ as $r$ goes from $1$ to $n$'').\\[0.3cm]
The sum $\sum_{k=p}^n f(n+1) - f(n)$ is called a telescoping sum.\\[0.3cm]
The pi symbol is used in the same way as the sigma symbol described above, except that succeeding terms are multiplied instead of added:\\[0.3cm]
$\lim_{x \to x_0} f(x) = \ell$ ``the number $\ell$ \textbf{is the limit of $f(x)$ as $x$ approaches $x_0$}''.\\[0.15cm]
$\lim_{x \to x_0^+} f(x) = \ell_1$ ``the number $\ell_1$ \textbf{is the right-hand limit of $f$ at $x_0$}''.\\[0.15cm]
$\lim_{x \to x_0^-} f(x) = \ell_2$, ``the number $\ell_2$ \textbf{is the left-hand limit of $f$ at $x_0$}''.
\vspace{0.15cm}
\end{minipage}%
}
\end{minipage}

\end{document}